\subsection{One Mumford factor and powers of the sixfold}\label{v7:ssec:mixed-single}

Fix a full-group fibre of the construction in Section~\ref{v7:sec:setup}.
Write $V=H^1(X,\QQ)$, $D=H^1(Z,\QQ)$ and $A_i=\operatorname{Sym}^2V_i$
after scalar extension to a common splitting field. The joint Hodge group
on $V\oplus D$ is the image of $G\times T$. Indeed its projection to the
adjoint group of $G$ is full because the fibre is full-group for $X$.
Its connected derived group is therefore the image of $G$. The projection
to the torus quotient is full by the cocharacter calculation in
Proposition~\ref{v7:prop:sixfold}. The finite central kernels do not affect
invariant tensors. In particular $V$ and $D$ have no common irreducible
constituent. There are no cross divisor classes between $X$ and $Z^m$.

\begin{proposition}\label{v7:prop:single-x-codim2}
For every $m\ge 1$, the rational codimension-two Hodge classes on
$X\times Z^m$ are spanned by products of divisor classes.
\end{proposition}

\begin{proof}
Let $D^{\oplus m}$ denote $H^1(Z^m,\QQ)$. The two weights of $T$ in each
block $D_i$ are inverse, and the characters from the three blocks are
independent over the splitting field. Thus $T$-invariant tensors in
$\bigwedge^q(D^{\oplus m})$ have even $q$. Only the K\"unneth bidegrees
$(4,0),(2,2),(0,4)$ can contribute. In degree two one has
\[
 \bigwedge^2V=\mathbf 1\oplus\bigoplus_{i<j}A_i\otimes A_j.
\]
In fact, of the three factors $V_i^{\otimes 2}$, either three or one
must be alternating to give the alternating part of $V^{\otimes 2}$.
The neutral part of $\bigwedge^2(D^{\oplus m})$ has only trivial and
single-adjoint constituents. A neutral wedge pair lies in one block,
where its two $T$-weights are inverse. Hence the only invariants in
K\"unneth bidegree $(2,2)$ are products of the polarization on $X$
with divisor classes on $Z^m$. In bidegree $(4,0)$ the invariants are
spanned by the square of the polarization on $X$ by the multiplicity
calculation of Proposition~\ref{v7:prop:square}. Bidegree $(0,4)$ is
generated by divisor products by Proposition~\ref{v7:prop:sixfold}.
The representation argument computes dimensions after scalar extension,
and the equality of the rational spaces follows by faithful descent.
\end{proof}

\begin{proposition}\label{v7:prop:single-x-codim3}
Let $m\ge1$. The quotient of $\Hdg^3(X\times Z^m)$ by the span of
products of three divisor classes has dimension $3m^4$ over $\QQ$.
It is supported in K\"unneth bidegree $(2,4)$. An algebraic class with
nonzero image in this quotient exists if and only if the full-group
mixed link $\xi$ is algebraic.
\end{proposition}

\begin{proof}
Torus neutrality again leaves only even K\"unneth bidegrees. In bidegree
$(6,0)$, hard Lefschetz identifies the invariants with those in $H^2(X)$,
so they are spanned by the cube of the polarization. Bidegree $(0,6)$
is generated by divisors by Proposition~\ref{v7:prop:sixfold}.
In bidegree $(4,2)$ only a trivial constituent of $\bigwedge^4V$ can
pair with $\bigwedge^2(D^{\oplus m})$. To check this, write $w_{abc}$
for the multiplicity of weight $(a,b,c)$ in $\bigwedge^4V$.
The weights of $V$ are the eight sign triples. The counts
$w_{200}=4$, $w_{220}=w_{202}=2$ and $w_{222}=1$ follow by choosing
four distinct sign triples, as in Proposition~\ref{v7:prop:square}.
The weight $w_{400}$ is one, since all four triples with first
coordinate $+1$ must be chosen, and $w_{420}$, $w_{402}$ and $w_{422}$
are zero for the same reason. Successive subtraction of raised
weights in the three coordinates gives the multiplicity of $A_1$ as
\[
 4-1-2-2+0+0+1-0=0.
\]
The other two adjoints have multiplicity zero by permutation. The
trivial constituent of $\bigwedge^4V$ is one-dimensional, so the
$(4,2)$ invariants are products of the square of the polarization and
divisors from $Z^m$.

In bidegree $(2,4)$ the trivial constituent of $\bigwedge^2V$ gives
only divisor products. Its complement is the primitive Hodge structure
$U_X$, with split representation $\bigoplus_{i<j}A_i\otimes A_j$.
To obtain $A_i\otimes A_j$ in neutral degree four of $Z^m$, one needs
two slots from block $i$ and two from block $j$. In each block, the
neutral pairs contribute one copy of $A_i$ for each choice of a plus
copy and a minus copy of $D_i$. There are $m^2$ such choices. Thus
each $A_i\otimes A_j$ has multiplicity $m^4$ in neutral degree four,
and the invariant mixed dimension is $3m^4$ after scalar extension.
This is also its rational dimension. Products of divisors in this
bidegree have $X$-component in the polarization line, because there
are no cross divisors. This identifies the quotient as claimed.

It remains to check the algebraic equivalence. In $H^4(Z,\QQ)$ let
$U_Z$ denote the rational summand with split constituents
$A_i\otimes A_j$ for $i<j$, each with multiplicity one. It descends
because these constituents form a single Galois orbit. Its Tate twist
$U_Z(1)$ is isomorphic to $U_X$ as a rational Hodge structure. More
explicitly, the two rational group representations become isomorphic
after scalar extension, and the space of equivariant maps commutes
with scalar extension. The rational structures are irreducible by
transitivity of Galois on their three distinct constituents, so a
nonzero rational equivariant map is an isomorphism. The common Hodge
circle gives the Hodge comparison after the indicated twist. Both
commutants are the cubic field $F$, obtained from the same three
embeddings. By Proposition~\ref{v7:prop:sixfold}, all Hodge maps
between cohomological summands of powers of $Z$, with Tate twists, are
algebraic. To use this assertion, identify such maps with Hodge classes
by Poincar\'e duality and use the algebraic degree projectors of an
abelian variety. In particular the projector and the $F$-action on
$U_Z$ are algebraic, and the $U_Z$-isotypic summand in $H^4(Z^m)$ can
be projected onto copies of $U_Z$ by algebraic correspondences.

Suppose a codimension-three algebraic class on $X\times Z^m$ has
nonzero quotient image. Its $(2,4)$ component defines a Hodge map
$H^2(X)\to H^4(Z^m)(1)$ using the Lefschetz pairing on $H^2(X)$.
This map is algebraic: multiply on the $X$-factor by the square of
its polarization and push forward. Project off the polarization
line on $X$ and then choose one of the algebraic maps to $U_Z(1)$
so that the resulting algebraic correspondence $f:U_X\to U_Z(1)$
is nonzero. It is an isomorphism. The algebraic adjoint and the
polynomial inverse of $f^\dagger f$ on $U_X$ transfer the whole
$F$-action, exactly as in Proposition~\ref{v7:prop:transfer}.
The pairings are nondegenerate on these summands: $U_X$ is primitive,
and the irreducible self-dual constituent $U_Z$ has multiplicity one
in its degree, so the invariant Lefschetz pairing restricts
nondegenerately. This supplies the primitive cubic action required in
Proposition~\ref{v7:prop:square}, and therefore all Hodge classes on
every power of $X$ are algebraic.

For clarity, one can recover the original mixed link from $f$.
In $U_Z^{\otimes 2}(3)$ the adjoint sum $T_Z$ has a nonzero Hodge
embedding. After splitting, for each $i$ choose a $j\ne i$ and use
the inclusion $A_i\hookrightarrow A_i\otimes A_i$ together with
the invariant pairing tensor in $A_j\otimes A_j$. Here the first
inclusion is dual to the nonzero Lie bracket on the adjoint module.
Thus there is a nonzero equivariant map after scalar extension.
Equivariance descends by extension of scalars on the space of maps,
and any nonzero rational map is injective since $T_Z$ is irreducible.
This map is algebraic by the Hodge conjecture on powers of $Z$.
Apply the algebraic inverse of $f^{\otimes 2}$ on its image to obtain
an embedding of $T_Z$ in $U_X^{\otimes2}(1)$. Its image is Hodge
isomorphic to $T_X$. Choose a Hodge map $U_X^{\otimes2}(1)\to T_X$
restricting to that isomorphism, using semisimplicity of polarized
rational Hodge structures. This last map is algebraic by the
now-established Hodge conjecture on powers of $X$. The composite gives
a nonzero algebraic map $T_Z\to T_X$, so its adjoint and the algebraic
$F$-action on $T_Z$ yield $\xi$ by Proposition~\ref{v7:prop:transfer}.

Conversely, assume that $\xi$ is algebraic. There is an algebraic
isomorphism from $T_X$ to $T_Z$ and an algebraic inverse on these
summands. After splitting, consider the two adjoint summands of
$T_X$ of colors $i$ and $j$, for $i\ne j$. Identify $V\otimes V$
with $\operatorname{End}(V)$ via the polarization. An element of
the $i$-th adjoint summand then acts as a trace-zero endomorphism
on $V_i$ and as identity on the other two factors. Products with
the $j$-th adjoint summand are self-adjoint for the alternating
form on $V$, since the two skew-adjoint endomorphisms commute.
They give the nonzero constituent $A_i\otimes A_j$ in
$\bigwedge^2 V$. Composition and the alternating projection are
algebraic operations, using the diagonal and the polarization,
as in Section~\ref{v7:sec:transfer}. On the $Z$ side the cup product
of the $i$-th and $j$-th adjoint classes of $H^2(Z)$ identifies
$A_i\otimes A_j$ with its copy in $U_Z$. An inverse on this
summand exists by semisimplicity and is algebraic by the Hodge
conjecture on powers of $Z$. Applying the two inverse links to
these factors and then composing on $V$ gives a nonzero algebraic
map $U_Z(1)\to U_X$. The construction uses rational maps throughout:
the maps to and from $U_Z$ can be restricted to the rational summand
of $T_Z^{\otimes2}$ with distinct colors, whose equivariant projector
is algebraic by the same result on $Z$. The map is nonzero on each
split constituent, hence its transpose, converted by algebraic
Lefschetz maps, gives a codimension-three algebraic class on
$X\times Z$ with nonzero quotient image. Pullback by a factor
projection $Z^m\to Z$ gives one on $X\times Z^m$.
\end{proof}
